\documentclass[11pt]{article}
\usepackage[T1]{fontenc}
\usepackage[utf8]{inputenc}
\usepackage{lmodern}
\usepackage[margin=1in]{geometry}
\usepackage{amsmath,amssymb,amsthm,mathtools}
\usepackage{enumitem,booktabs,tabularx,array}
\usepackage{microtype}
\usepackage[hidelinks]{hyperref}
\usepackage{tikz-cd}
\newtheorem{definition}{Definition}[section]
\newtheorem{target}[definition]{Proposed target}
\newtheorem{test}[definition]{Required test}
\theoremstyle{remark}
\newtheorem{remark}[definition]{Remark}
\newcommand{\Obs}{\operatorname{Obs}}
\newcommand{\Real}{\operatorname{Real}}
\newcommand{\Cost}{\operatorname{Cost}}
\newcommand{\im}{\operatorname{im}}
\newcommand{\rank}{\operatorname{rank}}
\newcommand{\conv}{\operatorname{conv}}
\newcommand{\spanlin}{\operatorname{span}}
\newcommand{\RR}{\mathbb R}
\newcommand{\NN}{\mathbb N}
\newcommand{\calE}{\mathcal E}
\newcommand{\calN}{\mathcal N}
\newcommand{\calR}{\mathcal R}
\newcommand{\calB}{\mathcal B}
\newcommand{\calC}{\mathcal C}
\newcommand{\calW}{\mathcal W}
\newcommand{\calU}{\mathcal U}
\newcommand{\calP}{\mathcal P}
\newcommand{\one}{\mathbf 1}
\newcommand{\Id}{\mathrm{id}}
\setlist{nosep,leftmargin=*}
\setlength{\parskip}{3pt}
\title{A research proposal for categorical phenomenology}
\author{Discussion draft}
\date{3 November 2026}
\begin{document}
\maketitle
\begin{abstract}
The central question is how structure can emerge from observation. Repeated response patterns can generate distinctions and concepts; repeated continuation patterns can generate states and dynamics; dependencies and factorization patterns can suggest operations, latent carriers, interfaces, and decompositions. What, if anything, is the common mathematical mechanism behind these passages? Which features are derived from the observations, which follow from a supplied mode of inference, and which remain choices?

This proposal seeks a common emergence and comparison theorem, or a rigorous classification of why several mechanisms are necessary. It starts with observational presentations rather than a preferred latent state space or component architecture. Its decisive output must actually construct structural presentations and relate their operations and possible decompositions, not merely organize already supplied models into a category. The existing generated-image, predictive, positive-realization, and two-sided cylinder results are substantial baselines. Native realization, refinement laws, composition, resources, and feasible acquisition then test what the emerged structure warrants and how it can be used.

This is a discussion proposal, not an adopted contract or a claim that the new theorem is known. It neither promises ontology discovery from an arbitrary experience stream nor treats a single motivating example as the scope of the theory. The aim is mathematical understanding across different ways in which patterns of observation become structure.
\end{abstract}
\tableofcontents

\section{The question and the proposed direction}
\label{sec:question}

``Turning observations into structures or models'' is the starting idea. A process law can generate predictive states even if no hidden state names were supplied. An incidence relation can generate a concept lattice even if no concept names were supplied. Observable dependencies can generate a linear dimension, while a factorization may introduce latent carriers or suggest a decomposition. These are forms of structural emergence in exact mathematics, not only statistical learning or discovery algorithms.

The central problem is to explain these passages and their relations. It begins before asking whether an existing representation is sufficient or preserved under a change. Such sufficiency and preservation questions are interesting and have substantial prior art in refinement, abstraction, and realization theory. They belong in the programme as tests and transformation laws for structure that has emerged. They should not define the whole programme.

The working aim is a \emph{common emergence and comparison theorem}: a construction, or a tightly related family of constructions, in which observational regularities yield structural presentations and the mathematics distinguishes identification, generation, invention, and choice. A presentation can contain derived states or objects, relations among them, operations induced by observational composition, and candidate interfaces or decompositions. The theorem should explain why different reconstruction mechanisms agree where they do, why they differ elsewhere, and how those differences affect native realization and use.

There is a plausible reason to regard this as a fundamental problem form. A forward observational semantics induces indistinguishability. Asking what can explain those observations raises existence, ambiguity, comparison, and choice, while repeated patterns can supply new structural objects. This is analogous in logical role to obstruction questions arising from existence problems. But the inverse image of a map is available everywhere. Defining a fibre does not establish a distinctive theory, and the fundamental-status claim remains something the programme must substantiate through reusable mathematics.

Three claims therefore remain separate:
\begin{enumerate}
\item \emph{Fundamental problem form}: observation-to-structure questions arise for structural reasons, with a common mathematical architecture worth understanding.
\item \emph{Canonical mathematics}: a particular emergence mechanism or comparison has a universal property relative to stated inputs.
\item \emph{New results}: a precise theorem adds something beyond existing concept analysis, realization, computational mechanics, categorical learning, abstraction, and composition theory.
\end{enumerate}
The first does not require a canonical optimal model. The second does not imply finite or physically available realization. The third requires independent comparison with the strongest sources.

Nor does emergence require that nothing be supplied. An experiment has outcomes, a record has a mode of comparison, and a continuation has an order. Closure under conjunction, randomization, linear combination, or composition can be a justified inference principle or a deliberate doctrine; it is not automatically discovered from a finite table. The task is to make the supplied/derived boundary exact, move it where mathematics permits, and reveal where competing structural choices remain.

The theory should transcend its examples in the way the contextad programme transcended its original motivation. The recent finite stochastic work provides stringent tests, but it does not determine the scope. The proposal allows a common construction if one can be proved, and a structured family with nontrivial comparison and impossibility theorems if that is the more faithful answer.

\section{What the existing work already establishes}
\label{sec:ledger}

The following ledger distinguishes the wider line of work from its recent finite stochastic testbed. ``Established here'' means proved or reviewed within the stated scope of the supplied notes, not formally verified or certified novel. Recovered exploratory claims have weaker standing than the six source-reviewed reports; their prerequisites must be rechecked before reuse.

\subsection{The wider mathematical baseline}

\begin{description}[style=nextline,leftmargin=0pt]
\item[Generated observable structure]
A supplied free/forgetful doctrine and response algebra yield a free extension of preparation profiles followed by image factorization. This includes the ordinary response quotient, meet completion, linear span, and convex hull. The generated image is initial among observable realizations and terminal among generated realizations in the appropriate categories. The Boolean/FCA comparison, preservation of generator relations, test-infomorphism comparisons, and explicit observation-change obstructions are already present \cite{Core,Recovered}.

\item[Predictive reconstruction]
For a fixed initial law and specified public adaptive protocols, histories can be identified by full future experimental laws. Filtering, reachable restriction, and observable minimization are distinct operations. Filtering has the stated adjoint direction in its supplied category. A finite hidden generator can have an infinite predictive quotient; hiding observations can require filtering afresh rather than induce a map on the old predictive states. Action translation and output hiding have different comparison laws \cite{Recovered}.

\item[Linear and positive realization]
Finite-rank Hankel reconstruction and its reachable/observable properties coexist with a nontrivial choice of invariant positive cone. Generated and test-positive cones give boundaries; polyhedral and simplicial cone criteria connect to classical finite and rank-sized positive realization. Linear exactness need not preserve the generated cone. Correlated preparations can defeat separately reconstructed local tensor spaces. These are substantive results and classical antecedents, not missing generic decompositions \cite{Recovered,Positive}.

\item[Two-sided cylinder normalization]
The completed cylinder and cross-doctrine work supplies coherent normalization spans in its adequate cylinder-factorization settings. It treats actual joint preparations/tests, fixed-port wiring, signed context variance, native comparison spans, and local exactness. The family-level and substitution coherence statements are already proved under their hypotheses. Some simultaneous changes require a span between apices rather than an ordinary map of spans. Descent from a free Isbell envelope may be only oplax and need not preserve the particular mixed pullback. Wire-type changes and trace need additional structure. This proposal must dovetail with that machinery, not replace it by a new untyped square \cite{Recovered,Garner}.
\end{description}

\subsection{The recent finite native testbed}

\begin{enumerate}
\item \textbf{Common generated images and native boundaries.} The core audit isolates the finite-simplex boundary and shows why an abstract convex image does not determine its latent stochastic cost. Distinct embeddings of an abstract square can have different minimum latent cardinalities. Separation of observed rows need not separate their algebraic completion \cite{Core}.
\item \textbf{Retaining a chosen model.} The refinement note separates coarse trace preservation from a strong labelled stochastic comparison to a particular chosen realization. Bounded synthesis reduces to existential real algebra. The two-stage class reduces exactly to Blackwell garbling and constrained nonnegative rank. Its partition family gives exact feasibility and costs, including an all-budget obstruction to retaining a late-randomness model \cite{Refinement}.
\item \textbf{A common native fragment.} Normalized Boolean and nonnegative-real factorizations share a finite two-stage construction. Independent support and convex criteria decide the specified unbounded feasibility problem; explicit formulas transport witnesses and give resource inequalities. Support need not reflect quantitative feasibility or optimal size \cite{Common}.
\item \textbf{Joint witnesses and reconstruction.} Fixed local factors, occurrence labels, and statewise continuation rows give a finite tensor polytope. Feasibility constructs an actual native witness, with a dual obstruction and a sharp constrained-rank resource law. The split reconstruction is initial in a \emph{fixed tensor fibre}, not a canonical optimizer across all tensors. The note also gives direct one-step transition-witness assembly without presupposing a global readout, and a simultaneous-cover test for committed deterministic intermediate diagrams \cite{Joint}.
\item \textbf{Fixed budgets.} The gated polynomial encoding is exactly equivalent to a fixed native state budget and its displayed quartic is coercive. Global zero, positive certified global minimum, stationary point, and small numerical residual have distinct meanings. The compositional gradient theorem supplies a field-assembly identity, not a global solver \cite{Budget}.
\item \textbf{Changing interfaces and optimization.} Lawful native transport gives
\[
 \kappa(D')\leq \min_{Z\in\calP(D)}\kappa_{D'}(F_*Z)\leq\kappa(D).
\]
Both inequalities can be strict for different reasons. Inherited value functions compose by classical infimization, while independently defined native target costs need not agree with them. Canonical optimum selection, smooth slack lifting, off-zero loss preservation, and forward gradient transport can fail, including in the actual quartic rather than an unrelated surrogate \cite{Interface}.
\end{enumerate}

These six reports mostly explore one finite stochastic architecture. Their accumulation is not evidence that arbitrary multistage realization, empirical discovery, or general phenomenology is finished. Conversely, the wider cylinder, Boolean, predictive, and positive results prevent treating the line of work as merely that architecture.

\subsection{What is still not supplied}

No current result derives an unrestricted scientific vocabulary from an arbitrary experience stream. None makes finite noisy data a complete exact behavioral specification, infers hidden-state intervention access from public traces, or promises a canonical minimum native model. The joint theorem does not establish arbitrary multistage construction with stochastic intermediate links, adaptive intervention, and independently chosen assembly routes. Formal compact or Isbell completions do not by themselves supply obtainable native systems. These are boundaries to investigate or retain, not defects to conceal.

\section{The relative data of a phenomenological problem}
\label{sec:data}

The following is proposed working language, not a definition imposed on every instance in advance. Its adequacy is itself tested by the examples.

\begin{definition}[Operational presentation]
An operational presentation $\mathfrak D$ specifies:
\begin{enumerate}
\item a typed repertoire $\Sigma$ of available preparations, observations, actions, and declared intervention operations, with the finite or effective descriptions actually available;
\item executable protocols $\calE_\Sigma$, including the permitted sequential, adaptive, parallel, and wiring operations, and their domains of definition;
\item an outcome semantics $\calB_\Sigma$ and an observation interpretation for each proposed native class $\calN_{\Sigma,d}$, where $d$ records a structural doctrine;
\item an observed subcollection $E\subseteq\calE_\Sigma$ and either exact responses $H_E$ or a specified uncertainty set $\calU_E$;
\item commitments $C$ to retain, such as a chosen realization, comparison direction, decomposition, conditional kernel, or intervention response;
\item a resource accounting and an access model for acquiring further responses and executing a realization.
\end{enumerate}
\end{definition}

The repertoire can be revised. It is not claimed to be the final ontology. A proposed new primitive must have a specified empirical interpretation, measurement or intervention procedure, or an explicitly conjectural status. A translation of old protocols into a richer repertoire records what old observations mean after revision. If no such interpretation is available, the new vocabulary is a candidate hypothesis rather than an experimentally identified structure.

For a native model $X$, write $\Obs_E(X)$ for its accessible responses. The consistent realization collection is schematically
\begin{equation}
 \Real_{d,C}(\mathfrak D)
 =\{X\in\calN_{\Sigma,d}:\Obs_E(X)\in\calU_E, X\text{ satisfies the stated }C\}.
 \label{eq:consistent}
\end{equation}
For exact data, take $\calU_E=\{H_E\}$. Morphisms must retain their concrete meaning, such as deterministic quotients, stochastic intertwiners, algebra homomorphisms, or the appropriate comparison spans. Equation~\eqref{eq:consistent} is bookkeeping. Every inverse problem has an analogous fibre; merely defining it is not a research result.

\subsection{Forced structure and chosen structure}

A structural predicate is \emph{forced relative to $\mathfrak D$} if it holds in every admissible consistent realization, interpreted up to the declared comparison or equivalence. It is refuted if none satisfy it and undetermined if both possibilities occur. This trichotomy presupposes a nonempty consistent class; an empty class is model failure, not vacuous evidence for every claim. A basis, positive cone, optimizer, coupling, decomposition, regularizer, or prior can be a useful choice without being forced.

A positive identification claim must name the invariant being identified. Isomorphic observable quotients, equivalent experimental predictions, equal state minima, the same retained transition mechanism, and isomorphic native models are different answers. A same-source choice of one realization does not prove that another route produces exactly the same output, or even an object in the same retained-witness fibre.

\subsection{Four quantifiers that cannot be interchanged}

For an observation change, one can ask for:
\begin{enumerate}
\item an abstract reconstruction and its comparison;
\item existence of \emph{some} native realization of the revised behavior;
\item a certificate relating \emph{specified} endpoint realizations;
\item a lift for \emph{every} admissible target realization, possibly within every budget.
\end{enumerate}
The fourth is much stronger than the second or third and is not the default demand. Useful forgetting need not lift every coarse choice. A proposed sufficiency theorem must state which quantifier it proves, rather than rejecting useful abstractions because unrestricted fixed-endpoint lifting fails.

\subsection{Resource semantics}

Use a declared ordered resource set, for example
\[
 \Cost(X)=(\text{latent states},\text{description size},\text{execution cost}),
\]
and keep experimental cost separately as a budget on protocols and samples. Addition or composition of costs requires a model; it is not inferred from category notation. State count is not automatically bits, table size, energy, time, or sampling complexity. External clocks, shared randomness, hidden correlations, preprocessing, and oracle calls must be charged or explicitly provided.

For a scalar native budget $b$ and tolerance $\varepsilon$, a basic profile is
\begin{equation}
 \mathscr B_{d,C}(\mathfrak D)
 =\{(b,\varepsilon):\exists X\in\calN_{\Sigma,d},\ C(X),\
 \Cost(X)\leq b,\ \exists H\in\calU_E:\ d_E(\Obs_E X,H)\leq\varepsilon\}.
 \label{eq:budget-profile}
\end{equation}
Here the metric or discrepancy and its relation to confidence must be specified. The displayed existential quantifier expresses \emph{some} compatible response; robust conclusions valid for \emph{every} member require a different quantifier. Exact feasibility, approximate feasibility, and robust feasibility must not share one symbol without this distinction. Retaining a resource spectrum or a set of witnesses is often appropriate when choosing an optimum coherently is impossible.

\section{Where structure can actually emerge}
\label{sec:emergence}

A useful theorem must identify structural output that was not installed in the input. Start, in a simple two-sided case, with a response pairing
\begin{equation}
 H:P\times T\longrightarrow\Omega,
 \label{eq:pairing}
\end{equation}
where $P$ and $T$ name available presentations or preparations and observations or tests. They need not describe an interactive system: a static incidence relation or a family of measured features is already an instance. Further composition or context data are included only when the observations support them or the experimental setup supplies them.

\subsection{Distinctions and states}

Equality of response profiles gives the observational relation
\[
 p\sim_Tp'\quad\Longleftrightarrow\quad H(p,t)=H(p',t)
 \text{ for every }t\in T.
\]
Its quotient contains distinctions made by these tests. In a temporal setting, take preparations to be accessible histories and tests to be continuations; equality of complete future laws generates predictive states. This exact construction assumes full behavioral semantics, an adequate exact oracle, or a justified ideal limit; a finite observed record alone does not supply those laws. The source of the state space is then a regularity of observed behavior rather than a chosen hidden carrier.

This is an identification relative to a language of tests. It does \emph{not} force every compatible native model to have literally identical underlying states for equivalent histories. Adding tests can split observational classes, while adding evidence within a fixed language can shrink a set of explanatory models. Those are different operations with different variance. A theorem must not conflate them.

\subsection{Relations and completed objects}

Incidence patterns can generate closed concepts or implications; linear dependence can generate a dimension and equations; admissible mixing can generate convex structure. A completion can introduce objects that were not individually observed. Its warrant may be mathematical closure under a supplied principle, rather than evidence that each formal object is an executable preparation.

The existing generated-image theorem already explains several such constructions. The new task is not to repeat it with more examples. It is to explain its relation to structures generated by continuation, latent factorization, and possible decomposition, including cases where their universal properties or native meanings are different.

\subsection{Operations induced by observed composition}

Suppose experimental contexts act on preparations and tests with an observed compatibility law such as
\begin{equation}
 H(c_*p,t)=H(p,c^*t).
 \label{eq:context-duality}
\end{equation}
One can ask whether $c_*$ descends to the emerged distinctions, whether its composites induce relations among the new operations, and whether those relations survive a chosen completion. The experimental context is supplied; its operation on a newly constructed state space and the relations it satisfies can be derived. Conversely, if no context action or compositional evidence is available, an invented multiplication on the state set must be labelled as additional structure.

A more ambitious question is whether patterns of contextual transformations determine a smaller operation language or reveal reusable action families. This is a discovery problem relative to the available contexts and admissible hypothesis language. It is not permission to claim all scientific operations follow from an unstructured sequence of outcomes.

\subsection{Latent carriers and decompositions}

A factorization pattern may introduce intermediate objects through which observed relations or laws pass. These objects can explain response dependencies or offer a compact implementation. Unlike a quotient or a closure, a smallest factorization need not be unique or characterized by the same universal property. Different factorization witnesses can carry different conditional mechanisms while agreeing on all present observations.

A decomposition is stronger than merely factoring a matrix through an arbitrary carrier. It posits reusable parts, interfaces, and an assembly operation. The structural question is whether patterns in the observations and available contexts warrant such parts and their reuse. At least one central test should construct candidate interfaces or factorizations without taking the desired component split as input. A declared class of admissible decompositions is allowed; supplying the actual factors, interface, and global witness and then proving they compose is a different, downstream question.

\subsection{A concrete supplied and derived ledger}

For every example and theorem, list:
\begin{enumerate}
\item the raw observational data and experimentally available manipulations;
\item the comparison of outcomes and inference principles supplied;
\item the structural objects, operations, relations, and interfaces constructed;
\item the invariant or universal property that warrants each construction;
\item the unresolved choices and the additional observations that could distinguish them;
\item the native objects or operations that the construction does and does not realize.
\end{enumerate}
This ledger should show genuine emergence even in exact, non-algorithmic mathematics. It also prevents a supplied latent category from quietly carrying all the structure the theorem claims to discover.

\section{The central target}
\label{sec:central}

\begin{target}[Common emergence and comparison]
\label{target:central}
Construct, for a stated but substantial class of observational presentations, structural presentations generated by their patterns of distinction, dependency, context action, and factorization. Relate at least the incidence/closure, temporal/predictive, and native-factorization mechanisms by a common construction or explicit comparison theorems, with source-level criteria explaining when their objects, operations, and candidate interfaces agree or necessarily differ. Prove a result about the structural output that is not already a generic image, fibre, or formal-completion theorem.
\end{target}

The first research deliverable must compare admissible choices of $\Gamma$, including what FCA, CALF, predictive reconstruction, and factorization methods already supply. An unrestricted hypothesis language can carry the answer in advance; one must show that the alleged output was not encoded as a hidden input. The central unknown is mathematical, not just terminological: what common operation on observational data explains why these different structures emerge? An answer might use two-sided completion, presentations by generators and relations, reconstruction from contexts, or a different categorical device. It might instead prove that the mechanisms divide into unavoidable types with a precise comparison calculus. Choosing the language first and fitting examples into it would not answer the question.

\subsection{A theorem shaped construction problem}

Write an observational presentation as $\mathfrak O=(P,T,H,\Gamma)$, where $\Gamma$ contains only the available contextual/compositional data and explicitly declared inference principles. Seek a construction
\begin{equation}
 \mathfrak E(\mathfrak O)
 =\bigl(\text{derived distinctions and objects},\
        \text{induced operations and relations},\
        \text{candidate structural explanations and comparisons}\bigr).
 \label{eq:emergence-output}
\end{equation}
The entries in~\eqref{eq:emergence-output} are obligations to make precise, not a finished definition. A category containing all possible explanations is too easy to form. A successful construction must produce usable invariants or presentations and prove why they arise from $\mathfrak O$.

The proposed theorem should have the following linked content:
\begin{enumerate}
\item \textbf{Generation.} Give an explicit data-to-structure construction. State precisely which objects or operation symbols are created, which equations are inferred, and which completion or factorization choices remain. The construction must not start with the desired hidden carrier or final component architecture.
\item \textbf{Comparison of mechanisms.} Recover the established incidence, predictive, and positive/latent constructions at their correct types and universal directions. Prove at least one source-level bridge which explains more than the fact that each output has a profile map. This may compare induced operation structures, admissible decompositions, or native realizability rather than identify the carriers.
\item \textbf{Structural emergence beyond a fixed quotient.} Derive candidate operations, relations, or interfaces from patterns in the observational/context data. At least one result must concern a decomposition or intermediate language not prespecified as the input factors. Its nonuniqueness or nonexistence may be part of the conclusion.
\item \textbf{Warrant and ambiguity.} Identify what is determined up to observational equivalence, what is determined only after a doctrine or structural use is selected, and what remains a choice among explanations. Provide independently checkable criteria or explicit invariants in the concrete classes, not only universal descriptions of the desired output.
\item \textbf{Use and revision.} Connect the generated presentation to actual native realizations and its downstream refinement, composition, resource, and acquisition questions. These test the construction; they do not replace its generative part.
\end{enumerate}

A shared pairing or a universal category can organize these demands but cannot satisfy them alone. Likewise, the existing free generated-image and cylinder normalization theorems already supply real common mathematics. A new result must locate an additional structural consequence or a precise classification boundary beyond that baseline.

\subsection{The decisive generative milestone}

\begin{target}[An observation to structure bridge]
\label{target:generative}
From a common observational presentation with available context composition, derive a structural presentation containing both a generated state/object system and its induced operations. Prove an independent criterion relating this presentation to at least two unlike emergence mechanisms, including one for which the native or factorization structure is not merely the same image construction. Use the criterion to construct or obstruct a reusable interface or decomposition whose factors were not supplied in advance.
\end{target}

A concrete first problem is to start with response patterns and a finite presentation of contextual transformations, rather than named hidden components. Determine which observational equivalences are stable under those contexts, which operations descend, and which compatible families of transformations and response dependencies admit a nontrivial factor presentation. In a deterministic/algebraic case, the output may use congruences, induced transformation algebras, or presentations by relations. In a stochastic/positive case, it must account for positivity, correlated observations, nonunique latent factors, and the difference between predictive and hidden state. Classical algebraic decomposition, concept analysis, automata reconstruction, positive realization, and domain completion are the first reduction targets.

The hoped-for theorem is not that a unique decomposition always exists. It should distinguish an identified structural invariant, several observationally equivalent factor presentations, and an obstruction to any factorization of the specified operational kind. If a gauge or equivalence class is the right answer, define it. If no finite native decomposition follows from a canonical behavioral object, expose that gap. A parametrized family must retain actual inference or synthesis difficulty, rather than become trivial after all parts are supplied.

For the first finite factor-presentation test, rule out identity factorizations and a one-part decomposition. Require at least two proper observable parts, each carrying distinctions used by an actual test or context, with a specified assembly that reconstructs the claimed joint responses and operation action. In a product-style test this can be witnessed by proper projections from the derived state system, joint separation by the projections, and verified action/readout equations. The projections, assignment of contexts to parts, and assembly witness are outputs to be found, not the desired factorization supplied in advance. An overlapping-interface version must state its corresponding nontriviality criterion. Verification uses the input response/context equations and a completeness assumption when extrapolating beyond them; matching a selected subtable is not enough. These conditions do not force uniqueness and do not guarantee native realizability. They make the decomposition question falsifiable.

This milestone cannot be satisfied by applying three known algorithms and placing their outputs next to one another. It requires a bridge theorem with a source-level criterion and a consequence for the generated structure. Conversely, a rigorous classification proving that no proposed common mechanism can have the demanded properties could be a major result. That outcome must establish the boundary broadly and retain the successful specific constructions; a single counterexample does not settle the programme.

\subsection{Strong emergence precedents}

Existing work already derives substantial structure rather than merely validating a supplied representation. Computational mechanics constructs predictive states and their transitions from process laws, while preserving a supplied alphabet, time organization, and predictive criterion \cite{ComputationalMechanics}. FCA derives lattice objects and operations from incidence; the categorical learning framework CALF derives an image hypothesis and a compatible algebra structure from an observational wrapper. Its Theorem 9 characterizes that structure by closedness and consistency under its factorization and endofunctor assumptions \cite{FCA,CALF}. These are direct emergence mechanisms.

Other precedents move the boundary further. Contextual grammar learning can construct contextual features and composition rules from positive strings and membership queries under finite-context and finite-kernel assumptions \cite{ContextualGrammar}. Causal macro-variable work constructs partitions and possible subsidiary decompositions with an explicit micro-level causal and intervention semantics \cite{MacroVariables}. Structural-form inference chooses among trees, rings, grids, and related forms within a graph grammar and probabilistic model \cite{StructuralForm}. Program-library learning constructs reusable operations within a supplied typed language and prior-guided search \cite{DreamCoder}.

These accomplishments are not interchangeable. Canonical structure relative to an observation interface, statistical identification under assumptions, and useful structure selected by a prior or compression criterion are distinct. The proposal should compare them without either dismissing noncanonical discovery or calling a chosen structure forced. Existing nonidentifiability results also show why a bare observational distribution need not determine a preferred latent factor ontology \cite{Disentanglement}. The possible added question is a precise common emergence/comparison law, including how operation signatures and interfaces are licensed by different observational syntax. That is a research question, not a cleared novelty claim.

\section{Refinement laws as a component of the programme}
\label{sec:refinement-component}

After a structure has emerged, one can ask what survives a richer test language, a changed doctrine, a new connection, or a replacement mechanism. This is refinement-law territory, with substantial prior art. The existing programme already contains strong comparison results, and abstract interpretation supplies soundness, completeness, and constructive domain repair; refinement mappings can require history and prophecy variables; interface theories give compositional substitutability under explicit compatibility conditions \cite{AbstractCompletion,RefinementMappings,InterfaceAutomata}.

Generalized strong preservation already connects observation languages, abstract domains, behavioral equivalences, and complete-domain construction; it is an especially strong common-theory control under its stated hypotheses \cite{StrongPreservation}. Predictive sufficiency can itself be a constructive emergence principle. The issue here is the narrower preservation question that starts from an already selected representation, not the word sufficiency.

The proposed refinement component should therefore answer a specific residual: how do the newly generated objects, operations, and candidate interfaces compare, and when can those comparisons be realized natively with the stated resources? The finite sufficiency test below is a demanding positive special case. Noninvertible comparisons, controlled forgetting, and set-valued explanations remain first-class outcomes. Universal same-budget lifting is not the definition of useful phenomenology.

\section{A multistage test for emerged interfaces}
\label{sec:multistage}

The recent native work gives a stringent downstream test of an emerged interface. This section does not replace Target~\ref{target:generative}; it asks whether such an interface can actually be used and revised.

\subsection{Input and requested output}

Take finite controlled systems with joint output and next-state kernels
\[
 M^i_{a,y}:S_i\longrightarrow\RR_{\geq0}^{S_i},
 \qquad \sum_y M^i_{a,y}\one=\one,
\]
finite supplied preparations, an explicitly declared finite operation/intervention repertoire, and a diagram of partial local observations and retained models. Include the prescribed deterministic or stochastic comparison maps with their directions. Require sequential state reuse through at least two successive nonterminal transitions: a state reached by the first operation can be acted on and observed by the next. When feedback is claimed, additionally allow a later action to depend on an earlier observed outcome. At least one intervention must replace or control a specified mechanism, rather than merely relabel an output. The permissible intervention semantics and access to it must be supplied; passive observations do not create that access.

A candidate joint implementation has its own finite carrier $V$, kernels $X_{a,z}$, preparations, and retained links $L_i:V\to\Delta S_i$ where stochastic links are intended. For a declared observation map $q_i$, a required retention equation has the concrete form
\begin{equation}
 \left(\sum_{z:q_i(z)=y}X_{a,z}\right)L_i
 =L_iM^i_{a,y},\qquad \pi_XL_i=\pi_i.
 \label{eq:multistage-retention}
\end{equation}
A stochastic output channel uses the corresponding weighted sum. Analogous equations are required for each declared intervention, with the exact operators it changes. These are statewise conditions, stronger than equality of the one initial trace law. More general native classes may require a different typed formulation; Equation~\eqref{eq:multistage-retention} is the first concrete target, not a definition of all phenomenology.

The input must not already include a full joint witness. It may include incompatible or only partially determined local pieces, and previously chosen intermediates can matter. The requested output is a construction or obstruction for a joint native realization and for at least two specified assembly routes which begin from those pieces. Rebracketing a single supplied full tensor is a regression case, not the new theorem.

\begin{target}[Native use of an emerged interface]
\label{target:multistage}
Find an explicitly described class of the preceding inputs, closed under a useful collection of sequential operations and mechanism interventions, with an independent finite sufficiency/compatibility criterion which:
\begin{enumerate}
\item constructs joint native witnesses from the local presentation, with no assumed global coupling or unknown global readout supplied as a premise;
\item retains the specified stochastic links and operational comparisons, or gives an exact criterion distinguishing when only weaker public behavior can be retained;
\item compares two independently built legal assembly routes by actual native morphisms or a justified typed replacement, including any genuinely additional compatibility condition;
\item proves a nontrivial native resource bound and a matching lower bound on a parametrized family whose difficulty is not eliminated by a free clock, a full-history state, or a supplied witness;
\item identifies how a finite observational sufficiency certificate is used in the construction and how it changes when a test, operation, or intervention is added or hidden.
\end{enumerate}
\end{target}

A theorem for all finite systems is not presumed. Candidate restrictions include a finite continuation basis, a controlled overlap pattern, a specified cone class, bounded interface width, or a verifiable stability/separation condition. Their meaning and testing cost must be stated. A chosen class must still contain genuine stochastic intermediate links, unknown joint choices, sequential reuse, and a nontrivial intervention or its precisely diagnosed failure. A restriction that reduces every case to one already solved row polytope fails the milestone.

\subsection{What counts as an independent criterion}

Acceptable possibilities include a finite relation/closure calculation that yields the necessary native witness, a dual certificate with a construction theorem, or a finite positivity and invariant-subspace condition whose sufficiency is proved for the declared class. A finite cone or rank certificate can be useful when the generators and bounds are supplied; existentially asking for an arbitrary invariant cone is not an effective unbounded criterion by itself. A criterion stated as ``there is a compatible joint realization'' or ``the operation preserves the relevant class'' must be replaced by an input-level test.

The proof must reduce standard cases to their established explanations: ordinary marginal extension, conditional products, supported coupling, Markov semipullbacks, multicausal gluing, and bounded real-algebraic synthesis. The additional theorem must identify exactly what survives after prescribed stochastic anchors, finite memory, or alternate assembly routes are imposed. Merely adding the word ``mechanism'' to a fixed conditional-marginal constraint is not a new result \cite{Joint,Refinement,Multicausal}.

\subsection{Negative outcome and repair}

If a proposed condition is insufficient, construct a minimal or structurally transparent counterexample with exact verification. Then state the failed conclusion and test a repair: enlarging the retained interface, retaining a set of joint witnesses rather than one optimum, allowing posterior states, adding a necessary probe, changing the comparison class, or explicitly relaxing a retention requirement. The repair must preserve a useful operational claim and have a proved positive theorem with resource accounting.

A counterexample can refute the proposed unrestricted statement. It cannot satisfy all the positive clauses of Target~\ref{target:multistage}. If no nontrivial repaired class is found, this milestone remains incomplete and the broad version of the proposal must be revised. A tiny example that isolates an obstruction is valuable; it is not a substitute for a family retaining substantial inference or synthesis difficulty.

\section{A positive sufficiency test within the refinement component}
\label{sec:finite-sufficiency}

A demanding positive instance makes the refinement component more falsifiable. It is deliberately stronger than ordinary forward forgetting, which remains useful when it fails.

Let $E\subseteq E^+$ be finite observation interfaces in a declared finite-dimensional or finite-presentation regime, with restriction $q$. Let $\mathcal G$ be specified generating operations and exposed interventions, with their allowed wiring rules. Fix a native budget cap $B$. For an exact full response specification $H$, write $\Real^+_b(H,C)$ for its native implementations with cost at most $b$, and $\Real_b(qH,qC)$ for the restricted problem, using the correctly translated commitments.

\begin{target}[Finite operational sufficiency with native lift]
\label{target:finite}
Derive finite input certificates, under stated structural assumptions, which imply all of the following for every $b\leq B$:
\begin{enumerate}
\item observations on $E$ determine the required $E^+$ predictions throughout the protocols generated by $\mathcal G$, not just on the originally measured preparations;
\item the budget feasibility statements agree,
\begin{equation}
 \Real^+_b(H,C)\ne\varnothing
 \quad\Longleftrightarrow\quad
 \Real_b(qH,qC)\ne\varnothing;
 \label{eq:budget-exactness}
\end{equation}
\item for the stated class of coarse implementations, the certificates construct same-budget native lifts and comparisons, with a composition law under the generated wiring and intervention operations;
\item a corresponding weaker certificate gives a controlled one-sided comparison, resource bound, or explicit missing-interface obligation when exactness fails.
\end{enumerate}
The hypotheses must be verifiable in the presentation, not a prior assumption of a reverse native functor, budget-preserving lifting, or coherence of the desired outputs. The theorem must say whether its native conclusion covers every coarse implementation, only specified endpoints, or only existence.
\end{target}

Equation~\eqref{eq:budget-exactness} alone is weaker than the lifting clause. For attained scalar integer costs, equality of minima is equivalent to agreement of feasibility at every budget; agreement only for $b\leq B$ can hold while both minima exceed $B$ and differ. A multidimensional resource frontier can carry more information than a scalar optimum. The purpose is to derive an operationally meaningful native lift from a finite sufficiency certificate, rather than prove equality by assuming lawful transports in both directions, which the existing interface note already does.

\subsection{A clean control lemma}

There is a simple finite prototype that should be proved and used as a control, not counted as completion. Suppose the \emph{full admissible latent-profile polytopes} $K,K'$ are supplied, and a linear observation map $q$ satisfies
\begin{equation}
 q(K)=K',\qquad
 \ker(q)\cap\spanlin(K-K)=\{0\}.
 \label{eq:polytope-iso}
\end{equation}
Finite rational vertex/facet and rank checks can certify these conditions for appropriate explicit presentations. Require the actual generating operations to satisfy their well-typed intertwining equations on sufficient generators, and require explicitly specified joint correlation carriers as well as the local polytopes. Then an affine inverse on the admissible profiles is available without requiring a stochastic inverse on the entire ambient simplex. One can ask how that inverse lifts finite latent factorizations and the declared operations while keeping the same hidden carrier.

This is stronger than injectivity on observed rows and potentially weaker than a globally invertible observation channel. It still supplies much of the needed structure in advance. The resulting conjugacy argument may be elementary. It must be compared with existing tomography, statistical sufficiency, automata minimization, and causal abstraction before any novelty claim.

\subsection{The hard bridge}

The substantive refinement extension is to derive usable certificates of the relevant full admissible structure from \emph{available observational and operational data}, instead of supplying it as $K$. In positive linear settings, the existing generated and test-positive cones suggest a candidate interval
\[
 C_-(E)\subseteq C\subseteq C_+(E).
\]
One must track the observation embedding, normalization, all continuation/intervention operators, and retained native links. Equality of suitable bounds can make a cone forced; a gap records genuine freedom. Neither local cone equality nor a basis of marginal tests determines arbitrary joint correlations. Nor does a polyhedral generated cone automatically describe every latent profile admitted by a native implementation.

A useful theorem would identify finite conditions, such as separation, positivity reflection on the relevant span, operational closure, and controlled correlation structure, which imply the native conclusion in Target~\ref{target:finite}. It must derive the native lift rather than conceal it in the conditions. Conversely, an obstruction should distinguish insufficient observation from a genuinely incompatible retained model or resource budget. At least one criterion should remain useful when the canonical predictive object is larger than the hidden implementation, and at least one should expose the failure of purely local observational adequacy after composition.

The same pattern should have an algebraic instance: preservation of observed generators is insufficient unless the relevant generated relations and operation uses are preserved. The common theorem should explain why relation preservation, predictive continuation closure, and positivity/realization conditions occupy comparable roles, without falsely identifying their native categories or adjunction directions.

\subsection{What would refute this proposed direction}

If every plausible finite certificate either encodes the entire sought native solution, reduces to an existing theorem without added content, or excludes the mandatory hard cases, then the proposed sufficiency theorem has not been obtained. A rigorous classification showing that the mechanisms cannot admit the proposed common form would be a serious theoretical outcome. It would support a structured family of theories with sharp bridges, rather than a universal reconstruction. Such a result must establish the obstruction at the stated level of generality; one easy counterexample does not settle the fundamental-status question.

\section{Comparison with the cylinder and context machinery}
\label{sec:cylinders}

The existing two-sided machinery is a mathematical dependency and a test of the proposal. In each supported instance, supply an interpretation
\[
 J_{\mathfrak D}:\Real(\mathfrak D)\longrightarrow\operatorname{CylReal}(\mathfrak D)
\]
into the appropriate cylinder-realization setting. This map must send the \emph{actual} preparations, tests, native comparisons, and declared operations to their two-sided data. It must record which native information is forgotten. A freely completed formal object is not silently substituted for its native preimage.

For an honest cylinder realization $X$, denote the established normalization span schematically by
\begin{equation}
 \nu_{\mathfrak D}(X)=
 \left[X\xleftarrow{\mathcal M}G_X
             \xrightarrow{\mathcal E}R_{\mathfrak D}\right].
 \label{eq:normalization-span}
\end{equation}
The actual classes and variance are those of the existing theorem. A supported change $F$ has its established cross-doctrine comparison
\begin{equation}
 \beta_F:
 F(R_{\mathfrak D})\xleftarrow{\mathcal M}G_F
                \xrightarrow{\mathcal E}R_{\mathfrak D'}.
 \label{eq:cylinder-beta}
\end{equation}
These are baselines, not conjectures of this proposal.

The new theory must give a specified comparison from its interpreted native construction to these spans. Where native forward transport $u_*$ exists, the first required square is
\[
\begin{tikzcd}[column sep=large,row sep=large]
\Real(\mathfrak D) \arrow[r,"u_*"] \arrow[d,"J_{\mathfrak D}"'] &
\Real(\mathfrak D') \arrow[d,"J_{\mathfrak D'}"] \\
\operatorname{CylReal}(\mathfrak D) \arrow[r,"F"] &
\operatorname{CylReal}(\mathfrak D').
\end{tikzcd}
\]
The claim may be strict commutation, a specified natural comparison, or an appropriate span-valued version; the theorem must choose and type it. It must then prove that normalizing this square produces the existing comparison~\eqref{eq:cylinder-beta}, up to a specified isomorphism or noninvertible comparison of spans. Equivalently, the native route and the established two-sided route must be related by an actual morphism or two-cell in the supported span calculus, not by the phrase ``compatible with cylinders.''

For consecutive changes, this comparison must respect the established compositor and its unit and associativity laws. For simultaneous signed changes in preparations and tests, it must use the surviving span-between-apices construction where an ordinary interchange map is impossible. For fixed-port wiring it must compare the actual joint cylinder; for wire-type change it must supply the extra transport and account for reversed dual variance. Feedback or trace requires the original definedness and preservation conditions. Monoidality alone does not prove trace preservation.

The decisive extra content is reflection toward native semantics. The two-sided image may exist although no admissible finite native realization does. A faithful statement must expose the resulting lifting obstruction or resource cost. Likewise, a functor preserving factorization classes can transport a comparison without preserving the relevant mixed pullback exactly. The proposed sufficiency certificate should tell when native conclusions can be recovered, rather than simply reasserting the free comparison theorem.

Contextads are a possible language for the dependence of operations on context, revision, or partial undo. They become appropriate only if an explicit construction in this problem satisfies their required structure and yields a result not already given by the cylinder machinery. An institution, double category, fibration, profunctor, or operadic doctrine is similarly a candidate tool. No one of these is selected as the answer before the native criteria and comparisons are understood.

\section{Acquiring a sufficient observational interface}
\label{sec:acquisition}

Exact emergence theorems can start from complete response laws or incidence data. Acquisition asks how much of the structural presentation can be obtained from feasible observations, and which further tests can resolve competing structures. This is a staged extension of the central target, not a claim that emergence means only learning or that every pure theorem must immediately deliver a practical learner.

\subsection{Separate the access models}

\begin{enumerate}
\item \textbf{Supplied exact presentation.} A finite generator, response oracle, or complete algebraic presentation is given. Exact symbolic conclusions are relative to that input. A finite table without a model-class promise is not a complete infinite-horizon specification.
\item \textbf{Exact query learning.} Membership, equivalence, counterexample, reset, or intervention oracles are available only when explicitly granted. Their computational and physical costs must be stated. Existing categorical query-learning theorems are the first reduction target, not new results to rename \cite{CALF,FunctorialLearning}.
\item \textbf{Finite sampling.} Protocols produce random data, potentially with adaptive choice and dependence. Specify resetability, mixing or trajectory assumptions, accessible actions, horizon, observation noise, and a confidence level. Confidence sets must be justified under that sampling law.
\item \textbf{Resource-bounded distinguishability.} Only probes within a time, memory, horizon, or sample budget are accessible. Distinct mathematical behaviors may then remain indistinguishable. This can be a legitimate target of a compositional abstraction theory, including a later computational or cryptographic instance, rather than an inconvenient approximation to unrestricted equivalence.
\end{enumerate}

\subsection{The acquisition target}

Given observational data and a stated class or inference promise, seek an adaptive procedure which produces candidate objects, relations, operations, or interfaces from Target~\ref{target:central}, with certificates appropriate to its identification claim. When alternatives remain, it should identify an executable discriminating experiment or report an unresolved identified set with a justified lower bound. For a specified downstream use, a sufficiency certificate is one possible output. Finite termination or a finite sample bound needs assumptions such as finite distinguishability, a bounded model class, excitation/observability, a positive separation margin, or stability. They must be stated and, where possible, independently tested. A finite rank upper bound learned only by fitting a finite table is not a verified promise about all future experiments.

The procedure should be driven by an unresolved distinction, relation, factorization, or failed certificate in the exact construction: a missing algebraic relation test, an unresolved continuation, a separating positive effect, a joint-correlation probe, or a specific intervention. This is where categorical structure could do work beyond packaging, by transporting a known adequate experiment across modules or localizing the reason an observational interface is insufficient. A proof must quantify the advantage or limitation actually obtained; no blanket efficiency advantage over existing approaches is assumed.

Exact native minimum size is generally unstable under perturbations. A sampled estimate should not be rounded into an exact rank or zero/nonzero support claim without a margin. A more defensible output is a confidence family of resource profiles,
\[
 \left\{\mathscr B_{d,C}(H):H\in\calU_E\right\},
\]
with lower and upper conclusions whose quantifiers are explicit. A theorem about a particular optimum requires additional separation from the relevant realization boundary. A posterior or regularizer can choose a model inside the remaining ambiguity; it does not convert that choice into identification.

\subsection{Approximate composition}

For a chosen operational discrepancy, seek an error law for the declared wiring and interventions. A schematic target is
\[
 d_W\bigl(\Obs(WX),\Obs(W\widehat X)\bigr)
 \leq \omega_W(\varepsilon_1,\ldots,\varepsilon_m;\text{declared stability data}),
\]
where each local approximation has its own certified error. The modulus must follow from actual system and wiring properties. For feedback, conditioning on rare events, or long horizons it may grow badly or fail uniformly. A generic triangle inequality is a baseline; the useful theorem must identify the structural hypotheses making the bound informative and construct examples showing their necessity.

A no-go lower bound should use two admissible models that are indistinguishable under the allowed probes but differ on the structural use. Passive agreement with different intervention responses, finite-horizon agreement with later divergence, and locally indistinguishable components with different joint behavior are mandatory controls. A later experiment is informative only if it is actually executable under the declared access model.

\section{Resource laws and computational methods}
\label{sec:resources}

Resource theory belongs in the programme because a structure can be adequately specified but infeasible to obtain or realize. The desired laws should distinguish:
\begin{enumerate}
\item the cost of the canonical abstract or predictive representation;
\item the minimum cost among native implementations of the observable behavior;
\item the extra cost, or impossibility, of retaining chosen mechanisms or intermediate witnesses;
\item the information and execution cost of using the interface in a composite;
\item the experiment and computation required to certify any of these claims.
\end{enumerate}
The finite hidden/infinite predictive example and the unbounded joint-correlation-memory family show that these distinctions are substantive.

Under a lawful witness map $t$, the inherited value
\[
 (t_!c)(z')=\inf\{c(z):t(z)=z'\}
\]
has a standard associative elimination law. That identity is not a new central theorem and does not make $t_!c$ equal to an independently defined target cost \cite{Interface,Variational,SteinSamuelson}. A native-sufficiency certificate can prove equality, a controlled gap, or impossibility for the latter where appropriate. It should also say when coarse freedom creates new witnesses rather than merely permits refactoring a transported one.

Optimization and compositional gradient methods are possible computational tools. The audited hypergraph result gives a correct assembly of local gradient contributions through linear variable-sharing maps; it does not guarantee existence, global convergence, native interpretation of off-zero iterates, or forward trajectory transport under observation erasure \cite{Budget,Interface,Hanks}. The exact quartic and its counterexamples remain regression tests for any proposed optimizer bridge. A new parameterization, metric, or constraint-preserving dynamics is a separate choice requiring its own equivalence and convergence claims.

The computational obligation is modest in one respect and strict in another. General polynomial-time synthesis is not presumed. Exact witness verification, exact lower-bound or infeasibility certificates when claimed, bit complexity, and clearly reported failure of numerical methods are required. A practical heuristic can be useful while an exact problem remains open, but it must not be substituted for a theorem about native feasibility.

\section{Generality through unlike mathematical instances}
\label{sec:generality}

The theory should transcend a motivating example in the way the contextad work did. Applications provide independent tests of the same emergence and comparison mechanism, not separate definitions of phenomenology or a promise to solve all applications at once.

\begin{enumerate}
\item \textbf{Algebraic and logical observations.} Boolean response profiles, FCA or another generated algebraic doctrine. The test is whether the theorem distinguishes observed-row separation from preservation of generated relations and warranted operation use. A merely stochastic reinterpretation of this example is insufficient.
\item \textbf{Controlled stochastic processes.} Predictive behavior, hidden native implementations, chosen statewise comparisons, sequential wiring, and exposed interventions. The multistage milestone is one demanding native test, including finite hidden models with infinite predictive representations.
\item \textbf{Linear and positive systems.} Finite-dimensional behavioral reconstruction, invariant positive cones, latent realizability, and active identification under declared control and stability assumptions. The theorem must distinguish linear sufficiency from positive and resource sufficiency. Classical linear system identification provides a substantial positive baseline \cite{ActiveID}.
\item \textbf{A materially different later application.} Select one after the central emergence mechanism is clearer: experimental physics with restricted instruments; transactional or robotic systems with discovered operational semantics and partial undo; scientific model composition; or computationally bounded observational equivalence. The selected case must contribute an independent failure mode or theorem, not decorative terminology.
\end{enumerate}

A quantum or physical example must use its own admissible instruments, positivity, and operational access, not ordinary probability with renamed states. A transactional example must specify which changes are available and what partial undo preserves. A resource-bounded computational example must state the adversary/probe class and the intended compositional guarantee. None licenses physical ontology or causal claims from an abstract universal property alone.

Before claiming generality, instantiate the same substantive theorem in at least the algebraic and stochastic/positive settings, and work through a finite-dimensional linear setting or explain precisely which additional analytic assumptions it requires. Examples may serve different roles: a small counterexample isolates a boundary, whereas a parametrized family or real benchmark tests whether inference difficulty remains. The proposal requires both roles and keeps them distinct.

\section{Mandatory stress tests}
\label{sec:stress}

Every proposed criterion and comparison should be run against the following suite. Existing answers are controls; reproducing them is not a new milestone.

\begin{enumerate}
\item \textbf{A structure must not be hidden in the input.} Re-run the supplied/derived ledger after changing the inference doctrine or hypothesis grammar. State which structural output changes and why. A pairing can be fixed while meet completion, linear span, and native factorization ask different questions.
\item \textbf{Observation alone need not identify operations.} Even separating profiles on a finite set do not specify a binary multiplication when no composite observations or algebraic doctrine is supplied. Such a negative control locates the input boundary; it is not the proposed new theorem.
\item \textbf{Emergence can be noncanonical.} Different latent decompositions can have the same observation law. A prior or compression score can select one without identifying it. Seek a structural invariant, an equivalence class, or a specific additional probe rather than silently treating one choice as forced.
\item \textbf{Completion creates new distinctions.} Boolean rows $10$ and $01$ remain distinct after retaining the first coordinate, while their meet/top completion loses distinctions. The criterion must not equate adequacy for observations with adequacy for generated structure.
\item \textbf{Native geometry is not the abstract hull.} Affinely isomorphic square response images can have different minimum latent costs in different output embeddings. The theorem must retain the admissible response embedding and latent-profile class.
\item \textbf{Predictive size and hidden size differ.} A two-state hidden bit model has infinitely many predictive states. Hiding output can destroy the map on previous predictive classes. A new filter or a set-valued comparison may be necessary.
\item \textbf{Joint preparation information is missing locally.} Individually rank-one behavior can have joint rank two under correlated preparation. The theorem must not reconstruct that joint behavior from local reductions alone.
\item \textbf{A chosen mechanism can block refinement.} Early/late randomness has equal coarse trace laws but different compatibility with a newly exposed signal under a specified statewise comparison. More hidden states do not automatically repair the anchor.
\item \textbf{Full support does not imply quantitative liftability.} The two-stage full-support controls share a Boolean possibility pattern but differ in stochastic feasibility and native minimum size. Support-level adequacy is not numerical sufficiency.
\item \textbf{One-step joint feasibility does not settle sequential assembly.} Independently chosen intermediate completions may have no simultaneous cover. A fixed global tensor makes rebracketing coherent but does not solve this stronger problem.
\item \textbf{Local memory does not bound joint memory.} Fixed one-state local factors can require unbounded joint correlation memory. All correlation carriers, phase clocks, and shared random resources must remain visible.
\item \textbf{Value equality does not select an optimum.} Symmetry can prevent an equivariant minimum-state choice. A set of optimal objects or a resource correspondence may be the correct output.
\item \textbf{Feasibility transport does not transport the solver.} Actual quartic zero sets, residual norms, gradients, and global minimum values behave differently under interface changes. A new solver theorem must confront the exact counterexamples.
\item \textbf{Finite information leaves alternatives.} Models can agree on all tested protocols and disagree after an untested continuation or intervention. The claimed acquisition result must use its class promise and access model here.
\item \textbf{Feedback can amplify error.} An approximate compositional claim must display stability or other control, or report a rigorous loss of a uniform bound.
\end{enumerate}

A criterion need not make every test positive. It must return the right kind of negative or noninvertible answer without concealing the loss by changing the native problem. If all difficult tests are excluded by its assumptions, the claimed generality must be reduced accordingly.

\section{Stages and decision gates}
\label{sec:gates}

These are proposed research stages, not accepted obligations or a schedule. The primary dependency is
\[
 \text{source and mechanism audit}
 \ \longrightarrow\ \text{observation to structure bridge}
 \ \longrightarrow\ \text{multistage and cylinder comparison}
 \ \longrightarrow\ \text{resource and acquisition extensions}.
\]
A failure early in this chain cannot be compensated for by quickly producing later manuscripts.

\begin{description}[style=nextline,leftmargin=0pt]
\item[Gate A \quad Locate the emergence question precisely]
Produce the supplied/derived ledger for incidence/FCA, predictive states, CALF and functorial automata, positive and latent realization, grammar/feature discovery, and at least one structural-form or operation-vocabulary learner. Compare the admissible $\Gamma$ choices rather than fixing a rich enough language to encode every answer. Audit the strongest sources independently. \emph{Pass}: a precise emergence/comparison claim with an added obligation not supplied by the audited theorems is stated; novelty remains provisional. \emph{Fail}: the claim is only a generic fibre, image, wrapper, or a known algorithm renamed. Reformulate before building a framework.

\item[Gate B \quad Prove the observation to structure bridge]
Complete Target~\ref{target:generative}, with explicit generation of objects and operations and a criterion for an unsupplied candidate interface or decomposition. Relate at least two unlike mechanisms, and test the third. \emph{Pass}: a source-level criterion has a new structural consequence or a proved classification boundary; the result is more than putting familiar outputs in one category. \emph{Fail}: all operations/factors are prespecified, the construction just repeats the known image theorem, or the bridge is an arbitrary common embedding. A broad impossibility classification can justify a revised thesis; an easy counterexample cannot close this gate on its own.

\item[Gate C \quad Relate emerged structure to native use]
Prove the relevant refinement and native-realization laws, including the explicit cylinder comparison of Section~\ref{sec:cylinders}. Use Target~\ref{target:multistage} as a demanding stochastic test, not the definition of the whole programme. Include a positive native result and a sharp obstruction/repair boundary, with resource bounds. \emph{Pass}: the emerged structural information explains actual operational use and the loss or choice in native realization. \emph{Fail}: only deterministic semipullbacks, static marginal extension, or rebracketing a supplied full witness is established as the supposedly new conclusion.

\item[Gate D \quad Relate emergence to feasible observation]
Specify one exact-query and one finite-sample access regime, or justify concentrating first on one while leaving the other explicitly open. Prove an acquisition, identifiability, or lower-bound theorem linked to a generated object, relation, operation, or candidate interface from Gate B. \emph{Pass}: the structural target, executable probes, assumptions, confidence/termination, and resource costs are explicit. \emph{Fail}: an equivalence oracle is called finite data, a prior-selected structure is called forced, exact boundaries are inferred without a margin, or the proposed probe is inaccessible.

\item[Gate E \quad Test transfer and synthesize]
Use the same emergence/comparison theorem in unlike mathematical instances, including static and sequential data. Determine whether the evidence supports one construction, a structured family with common laws, or a rigorous impossibility classification. \emph{Pass}: a reader can identify what the common mathematics proves about the structural outputs beyond the separate known constructions. \emph{Fail}: only a shared diagram or changed labels connect them.
\end{description}

Gate B is the main generative commitment one would consider adopting after review. Gates C and D test its mathematical and empirical meaning. Neither a successful refinement theorem nor a polished implementation can compensate for failure to identify an emergence theorem. Acquisition can be staged, but cannot be omitted when claiming a method for feasible discovery from experience. A severe negative theorem can be a major result even without a universal repair. In that case the original positive target remains refuted or incomplete rather than being retroactively marked ``completed.''

\section{Proof discipline and deliverables}
\label{sec:discipline}

Each substantive result should be accompanied by:
\begin{enumerate}
\item a full statement of native objects, observations, operations, comparison maps, and all finiteness, positivity, reachability, separation, access, and resource assumptions;
\item a proof in which the independent criterion is separated from the existence conclusion, and a typed diagram showing what commutes or only compares;
\item a strongest-source audit identifying the exact predecessor theorem, the reduction to it, and the additional claim, with full text checked where a precise theorem attribution is made;
\item exact examples, obstruction certificates, and reproducible finite checks for the claimed formulas, including zero masses, unused states, empty fibres, and infinite-cost cases;
\item an independent mathematical review and a separate source/provenance review; neither finite testing nor independent review is called formal verification;
\item a status ledger distinguishing proved results, conditional conjectures, counterexamples, repaired theorems, computational evidence, and unresolved tasks.
\end{enumerate}

The recent reports provide a good control standard. Their exact checks and independent reviews support the stated finite formulas; their source audits explicitly avoid novelty clearance. The fixed-budget audit distinguishes coercivity and exact certificates from optimizer success. The interface audit checks the actual residual family and variance, rather than substituting a convenient different objective. These standards should carry forward.

Suggested outputs are an emergence/comparison theorem manuscript, a chapter relating native realization and refinement to that construction, a source-and-boundary ledger, and a small exact verification suite. An acquisition paper follows only when its access and statistical assumptions are adequate. The deliverables are mathematical results, not a count of manuscripts or definitions. No deadline, publication commitment, implementation commitment, or adoption of this proposal is implied by this draft.

This deliverable is intentionally uncompiled \LaTeX{} source. Static checks cover source structure and cross-references; no PDF layout or compilation is claimed. The proposed new theorems remain research targets, and this draft does not adopt them as a contract.

\section{What the proposal would let us conclude}
\label{sec:conclusion}

If successful, the programme would explain how a structure arises from observational patterns before it is treated as a representation to preserve. It would relate generated distinctions, concepts, predictive states, induced operations, and latent or decompositional explanations, while distinguishing their warrants and choices. Its central contribution would explain how patterns of observation generate objects, operations, and candidate interfaces, and why different emergence mechanisms agree or divide. Native realization, refinement, and resources would establish what those structures can support.

If the strongest unification fails, a classification of the unavoidable divisions could still establish a fundamental architecture for these problems: what observations force, what realization adds, where choice enters, and which changes have only one-sided or resource-sensitive comparisons. The failure must be proved at the appropriate level, not inferred from one inconvenient example.

The proposed next decision is whether this central target and its decisive positive tests capture the desired research direction. The supplied results already justify taking the question seriously. They do not yet establish that categorical phenomenology is one canonical construction, a new mathematical discipline, or a universal method of discovering primitives. The proposal is designed to find out which stronger claim can actually be defended.

\appendix
\section{Assumptions that must remain visible}
\label{app:assumptions}

\begin{enumerate}
\item A repertoire of measurements/actions is supplied at each stage. Revision is allowed, but a new primitive needs an interpretation or an explicitly conjectural status.
\item Algebraic closure and native admissibility are separate inputs. Formal joins, convex mixtures, tensor products, or compact completions need not be executable states or operations.
\item A named native comparison has a direction and strength. Public traces, conditional kernels, statewise simulation, deterministic projection, and intervention preservation are not interchangeable.
\item A finite table is not complete behavior without a finite-presentation or model-class assumption. A finite-dimensional fit is not an independently established dimension bound.
\item Declared hidden commitments are not automatically empirically identified. A specified statewise kernel can be retained as a design commitment while remaining observationally underdetermined.
\item Local observational adequacy does not give joint correlation adequacy. Allowed correlations, shared randomness, preparation scopes, and wire interfaces are explicit.
\item Unbounded native existence, fixed-budget feasibility, exact optimum, approximate fit, and optimizer convergence are separate questions.
\item Budget preservation for every coarse model is stronger than existence, equality of optimum values, or comparison of chosen endpoints.
\item Noisy support, rank, and exact resource boundaries require margins or set-valued conclusions. Conditioning and feedback require their own error control.
\item Every enlargement of a model class or weakening of commitments changes the problem unless a comparison theorem proves the intended original use survives.
\end{enumerate}

\section{Source register and scope of attribution}
\label{app:sources}

The six local reports and their review files are the primary record of the recent work. The recovered ledger identifies the earlier Boolean/stochastic, predictive, positive-Hankel, cylinder, and cross-doctrine checkpoints. The draft uses those sources to set the baseline; it does not claim to have independently reproved every recovered exploratory theorem.

The external references below are deliberately selective. The emergence precedents in Section~\ref{sec:central} control claims of structural, feature, and operation discovery; Garner controls the generic cylinder/Isbell claims; functorial automata and categorical learning control generic reconstruction and exact acquisition; statistical comparison and causal abstraction control the experimental and intervention language; positive realization and multicausal coupling control native and temporal claims; variational analysis and the gradient paper control optimization. Their juxtaposition is not evidence that the proposed bridge is absent from the literature. Before a novelty assertion, the precise target must be checked against the full strongest sources, including their hypotheses and later developments.

McCullagh's design-indexed statistical models, Lorenz--Tull's query-indexed causal abstractions, and double-operadic systems doctrines are especially important controls: changing experimental design, naturality of abstraction, and compositional interface syntax already have substantive theories. The proposed emergence theorem must coexist with these theories, and its refinement component must add a specific structural or native consequence rather than merely reproduce their languages \cite{McCullagh,LorenzTull,SystemsDoctrines}.

\begin{thebibliography}{99}
\bibitem{Core}
\emph{A common generated observable construction and its native boundaries}.
Programme report, 2 November 2026. Accompanied by source audit, recovered theorem ledger, exact checks, and review.

\bibitem{Refinement}
\emph{A focused reduction test for refinement of a chosen stochastic realization}. Programme report, 2 November 2026.
Finite controlled stochastic refinement, exact real-algebra reduction, two-stage garbling and partition criteria. Accompanied by source audit and review.

\bibitem{Common}
\emph{A common native construction for finite relational and stochastic refinements}.
Programme report, 2 November 2026. Accompanied by independent lifting criteria, source audit, stress tests, and review.

\bibitem{Joint}
\emph{Joint native realization with fixed local mechanisms}.
Programme report, 3 November 2026. Finite tensor and row-polytope assembly, constrained-rank resources, witness-fibre initiality, and transport. Accompanied by source audit and review.

\bibitem{Budget}
\emph{Fixed budget native realization and compositional gradient assembly}.
Programme report, 3 November 2026. Accompanied by exact residual checks, source audit, and mathematical review.

\bibitem{Interface}
\emph{Interface changes and fixed budget optimization}.
Programme report, 3 November 2026. Accompanied by exact examples, source audit, and independent review.

\bibitem{Recovered}
Recovered theorem ledger, 2 November 2026.
Records the earlier Boolean and stochastic discussions, predictive construction, positive Hankel construction, completed cylinder/wiring checkpoint, and completed cross-doctrine checkpoint. Bibliographic and proof status varies by source.

\bibitem{ComputationalMechanics}
Cosma Rohilla Shalizi and James P. Crutchfield.
\emph{Computational Mechanics: Pattern and Prediction, Structure and Simplicity}.
Journal of Statistical Physics 104, 817--879, 2001.
\url{https://arxiv.org/abs/cond-mat/9907176}.

\bibitem{FCA}
Rudolf Wille. \emph{Restructuring Lattice Theory: An Approach Based on Hierarchies of Concepts}. 1982.
\url{https://doi.org/10.1007/978-94-009-7798-3_15}.

\bibitem{ContextualGrammar}
Alexander Clark, R\'emi Eyraud, and Amaury Habrard.
\emph{Using Contextual Representations to Efficiently Learn Context-Free Languages}.
Journal of Machine Learning Research 11, 2707--2744, 2010.
\url{https://jmlr.org/papers/v11/clark10a.html}.
Its identification-in-the-limit and update/query-work bounds do not give unrestricted finite-sample exact recovery.

\bibitem{MacroVariables}
Krzysztof Chalupka, Frederick Eberhardt, and Pietro Perona.
\emph{Multi-Level Cause-Effect Systems}. AISTATS 2016.
\url{https://proceedings.mlr.press/v51/chalupka16.html}.
Used for explicit variable/decomposition construction with supplied causal and intervention assumptions, not an unchecked minimal-statistical-sufficiency claim.

\bibitem{StructuralForm}
Charles Kemp and Joshua B. Tenenbaum.
\emph{The discovery of structural form}. PNAS 105, 10687--10692, 2008.
\url{https://pmc.ncbi.nlm.nih.gov/articles/PMC2492756/}.

\bibitem{DreamCoder}
Kevin Ellis and collaborators.
\emph{DreamCoder: Bootstrapping Inductive Program Synthesis with Wake-Sleep Library Learning}. PLDI 2021.
\url{https://doi.org/10.1145/3453483.3454080}.

\bibitem{Disentanglement}
Francesco Locatello and collaborators.
\emph{Challenging Common Assumptions in the Unsupervised Learning of Disentangled Representations}. ICML 2019.
\url{https://proceedings.mlr.press/v97/locatello19a.html}.
The stated nonidentifiability result is model-relative, not an impossibility theorem for all structural discovery.

\bibitem{AbstractCompletion}
Roberto Giacobazzi, Francesco Ranzato, and Francesca Scozzari.
\emph{Complete Abstract Interpretations made Constructive}. MFCS 1998.
\url{https://www.math.unipd.it/~ranzato/papers/mfcs98.pdf}.
Constructive completeness repair under its continuity assumptions.

\bibitem{StrongPreservation}
Francesco Ranzato and Francesco Tapparo.
\emph{Generalized Strong Preservation by Abstract Interpretation}.
Journal of Logic and Computation, 2007.
\url{https://arxiv.org/abs/cs/0401016}.
Forward and backward completeness have different equations and must not be conflated.

\bibitem{RefinementMappings}
Mart\'in Abadi and Leslie Lamport.
\emph{The Existence of Refinement Mappings}.
\url{https://lamport.azurewebsites.net/pubs/abadi-existence.pdf}.
Its existence theorem uses stated machine-closure, continuity, and nondeterminism assumptions and can add history and prophecy variables.

\bibitem{InterfaceAutomata}
Luca de Alfaro and Thomas A. Henzinger.
\emph{Interface Automata}. 2001.
\url{https://web.cs.wpi.edu/~heineman/html/teaching_/CS562/p109-de_alfaro.pdf}.
Compositional substitution requires its explicit compatibility conditions.

\bibitem{Garner}
Richard Garner. \emph{The Isbell monad}.
\url{https://arxiv.org/abs/1410.7108}.
Cylinder-factorization pseudoalgebras and lax/colax comparisons; not a minimum-native-realization theorem.

\bibitem{Automata}
Thomas Colcombet and Daniela Petri\c{s}an.
\emph{Automata Minimization: a Functorial Approach}.
\url{https://arxiv.org/abs/1711.03063}.

\bibitem{CALF}
Gerco van Heerdt, Matteo Sammartino, and Alexandra Silva.
\emph{CALF: Categorical Automata Learning Framework}.
\url{https://arxiv.org/abs/1704.05676}.
See also the algebraic/tree-automata development at
\url{https://arxiv.org/abs/2001.05786}.

\bibitem{FunctorialLearning}
Thomas Colcombet, Daniela Petri\c{s}an, and Riccardo Stabile.
\emph{Learning automata and transducers: a categorical approach}.
\url{https://arxiv.org/abs/2010.13675}.
Exact-query assumptions and termination hypotheses must be retained.

\bibitem{BlackwellCategory}
Tobias Fritz, Tom\'a\v{s} Gonda, Paolo Perrone, and Eigil Fjeldgren Rischel.
\emph{Representable Markov Categories and Comparison of Statistical Experiments in Categorical Probability}.
Theoretical Computer Science 961, 113896, 2023.
\url{https://arxiv.org/abs/2010.07416}.

\bibitem{Positive}
Mike Boyle and Karl Petersen.
\emph{Hidden Markov Processes in the Context of Symbolic Dynamics}.
Author manuscript, including the positive-realization account and Heller criterion.
\url{https://math.umd.edu/~mboyle/papers/BoylePetersen07-10-2009.pdf}.

\bibitem{Multicausal}
Beatrice Acciaio, Daniel Kr\v{s}ek, and Gudmund Pammer.
\emph{Multicausal transport: barycenters and dynamic matching}.
SIAM Journal on Financial Mathematics 16(3), 1104--1138, 2025.
\url{https://arxiv.org/abs/2401.12748}.
Successive kernel and gluing results are prior art; latent-cardinality claims require their own proof.

\bibitem{McCullagh}
Peter McCullagh. \emph{What is a statistical model?} 2002.
\url{https://www.stat.uchicago.edu/~pmcc/pubs/AOS023.pdf}.

\bibitem{LorenzTull}
Robin Lorenz and Sean Tull. \emph{Causal and Compositional Abstraction}, 2026.
\url{https://arxiv.org/abs/2602.16612}.
Query-indexed and mechanism-level causal abstraction, including interventions; precise attribution is subject to the target's source audit.

\bibitem{SystemsDoctrines}
Sophie Libkind and David Jaz Myers. \emph{Towards a double operadic theory of systems}, 2025.
\url{https://arxiv.org/abs/2505.18329}.
Used as a compositional-syntax control, not as an empirical identification theorem.

\bibitem{ActiveID}
Andrew Wagenmaker and Kevin Jamieson.
\emph{Active Learning for Identification of Linear Dynamical Systems}.
COLT 2020.
\url{https://proceedings.mlr.press/v125/wagenmaker20a.html}.

\bibitem{Variational}
R. Tyrrell Rockafellar and Roger J.-B. Wets.
\emph{Variational Analysis}. Springer, 1998.
\url{https://sites.math.washington.edu/~rtr/papers/rtr169-VarAnalysis-RockWets.pdf}.
Equation 1(17) is the pertinent value-function/infimization construction.

\bibitem{SteinSamuelson}
Dario Stein and Richard Samuelson.
\emph{Towards a Compositional Framework for Convex Analysis (with Applications to Probability Theory)}.
FoSSaCS 2024.
\url{https://arxiv.org/abs/2312.02291}.
The native rank and quartic objectives need not satisfy its convexity hypotheses.

\bibitem{Hanks}
Tyler Hanks, James Fairbanks, and Matthew Klawonn.
\emph{Generalized Gradient Descent is a Hypergraph Functor}.
EPTCS 429, 217--233, 2025.
\url{https://arxiv.org/abs/2403.19845}.
Theorems 3.4--3.5 and Corollary 3.5.1 concern field assembly and gradient-flow functoriality, with their additional hypotheses.
\end{thebibliography}
\end{document}
